\chapter{Arthritis in Children (Juvenile Arthritis)}
\index{Arthritis in Children (Juvenile Arthritis)}

\section*{Definition and Overview}
Juvenile arthritis, formally called \textbf{juvenile idiopathic arthritis} (JIA), describes a group of chronic inflammatory joint conditions that begin before the age of 16 and persist for at least six weeks. The term idiopathic reflects the fact that the cause is not yet fully understood. JIA is an autoimmune or autoinflammatory disorder in which the immune system attacks the lining of the joints, causing inflammation, swelling, and pain. It is the most common chronic rheumatic disease of childhood and is a distinct group of diseases, not simply adult arthritis occurring in a child. Several subtypes exist with different patterns of joint involvement, associated features, and long-term outlooks.

\section*{Epidemiology}
Juvenile idiopathic arthritis affects approximately one to two children in every thousand. It is slightly more common in girls and can begin at any age, with peaks between one and three years and around puberty. The most common subtype is oligoarticular JIA, which affects four or fewer joints and is seen most often in young girls. The condition occurs worldwide and across all populations. Importantly, it is not caused by anything the family did or did not do, and it is not contagious.

\section*{Aetiology and Pathophysiology}
In JIA the immune system, which normally defends the body against infection, mistakenly attacks the thin membrane lining the joints, known as the synovium. The exact trigger is unknown, but both inherited genetic susceptibility and environmental factors appear to contribute.

\begin{itemize}
    \item \textbf{Immune attack:} white blood cells invade the synovium, producing inflammation, swelling, and excess joint fluid
    \item \textbf{Genetic susceptibility:} certain gene patterns, particularly within the HLA system, make some children more likely to develop JIA
    \item \textbf{Environmental triggers:} an infection or other stressor may precipitate the immune response in a susceptible child
    \item \textbf{Subtypes:} oligoarticular, polyarticular, systemic, enthesitis-related, and psoriatic forms differ in their joints, associated features, and risks
    \item \textbf{Systemic JIA:} characterised by widespread inflammation causing quotidian fevers, a salmon-pink rash, and enlarged lymph glands
    \item \textbf{Associated eye disease:} uveitis, an inflammation within the eye, is a particular risk in some subtypes and can be silent
\end{itemize}

\section*{Clinical Features}
Children often do not complain of pain as clearly as adults, so swelling, limping, and changes in behaviour are important clues to the diagnosis.

\begin{itemize}
    \item Swelling, warmth, or fullness of one or more joints, most often the knee, ankle, wrist, or fingers
    \item Morning stiffness that improves during the day, or stiffness after periods of rest
    \item A limp, or a child who refuses to use a limb or hand
    \item Pain that may be surprisingly mild compared with the visible swelling
    \item Irritability, tiredness, and reluctance to be physically active
    \item In systemic JIA: daily high fevers, a salmon-pink rash, and swollen glands
    \item In some children: eye redness or light sensitivity from uveitis, which is frequently painless and may go unnoticed
    \item Growth problems, including a limb that grows unevenly because of persistent inflammation around the growth plate
\end{itemize}

\section*{Differential Diagnosis}
Many other childhood conditions cause joint or limb symptoms and must be considered before diagnosing JIA.

\begin{itemize}
    \item Transient (post-viral) synovitis, a short-lived inflammation of the hip or knee after a viral illness
    \item Septic arthritis, an emergency bacterial infection of a joint
    \item Osteomyelitis, an infection of the bone
    \item Leukaemia and other cancers, which can present with bone and joint pain
    \item Reactive arthritis, occurring after a gastrointestinal or genitourinary infection
    \item Growing pains, which typically affect the legs at night and settle by morning
    \item Injury or a minor undisplaced fracture
    \item Connective-tissue diseases such as systemic lupus erythematosus and juvenile dermatomyositis
    \item Haemophilia with bleeding into a joint, and sickle cell disease crisis
\end{itemize}

\section*{When to Seek Medical Care}
See a doctor if a child has joint swelling or stiffness lasting more than a few weeks, a persistent limp, or unexplained fevers with a rash. A child diagnosed with JIA should also have regular eye examinations, because uveitis can develop silently and cause permanent sight loss if missed.

\textbf{Contact your local emergency services (or go to an emergency department) for:}
\begin{itemize}
    \item A single hot, red, extremely painful joint, especially with fever, which may indicate septic arthritis
    \item A child who refuses to bear weight or to use a limb
    \item High fever with a rash and severe joint pain, or any other signs of serious illness
\end{itemize}

\section*{Diagnosis and Investigations}
There is no single diagnostic test for JIA. The diagnosis is clinical, defined as inflammation of one or more joints lasting at least six weeks in a child under 16, after other causes have been excluded.

\begin{itemize}
    \item \textbf{History and examination:} the pattern and duration of joint involvement, fevers, rash, eye symptoms, and family history
    \item \textbf{Blood tests:} inflammatory markers (ESR and C-reactive protein), full blood count, antinuclear antibodies (ANA), and occasionally rheumatoid factor
    \item \textbf{Imaging:} X-rays, ultrasound, or MRI of affected joints to assess synovitis, effusions, and any early damage
    \item \textbf{Joint fluid analysis:} fluid drawn from a swollen joint may be examined and cultured to exclude infection
    \item \textbf{Eye examination:} slit-lamp assessment by an ophthalmologist to detect uveitis
    \item \textbf{Referral:} to a paediatric rheumatologist for confirmation of the diagnosis and ongoing management
\end{itemize}

\section*{Management}
Treatment is delivered by a multidisciplinary team that includes a paediatric rheumatologist, physiotherapist, occupational therapist, and ophthalmologist. The goal is to control inflammation completely so that joints, growth, and vision are protected.

\begin{itemize}
    \item \textbf{Anti-inflammatory medicines:} non-steroidal anti-inflammatory drugs (NSAIDs) to control pain and swelling
    \item \textbf{Joint injections:} corticosteroid injections into affected joints for rapid relief of flares
    \item \textbf{Methotrexate:} a disease-modifying medicine that is the mainstay of treatment for many children
    \item \textbf{Biological medicines:} newer injectable drugs that block specific immune signals when conventional treatment is insufficient
    \item \textbf{Steroid tablets:} used for severe or systemic disease, kept as short and low-dose as possible
    \item \textbf{Physiotherapy and exercise:} keep joints moving, muscles strong, and physical function good
    \item \textbf{Ophthalmology care:} regular eye screening and prompt treatment of uveitis if it develops
    \item \textbf{School and family support:} education and practical support to keep the child learning, active, and socially connected
\end{itemize}

\section*{Complications}
\begin{itemize}
    \item Joint damage, contractures, and long-term disability if inflammation is not controlled
    \item Uveitis, which can cause permanent sight loss if not detected and treated
    \item Uneven growth, limb-length discrepancy, or short stature
    \item Thinning of the bones, partly from the disease itself and partly from corticosteroid use
    \item Chronic pain, fatigue, and reduced physical fitness
    \item Emotional and social difficulties, including missed schooling and low self-esteem
    \item Side effects of medicines, including increased infection risk with immunosuppressive drugs
\end{itemize}

\section*{Prognosis and Outcomes}
The outlook varies with the subtype of JIA. Many children experience long periods of remission, and a substantial proportion eventually outgrow the active disease. With modern treatment, most children lead full, active lives, attend school, and take part in sport, although some continue to require treatment into adult life and a minority develop permanent joint damage. Early diagnosis, prompt effective treatment, and regular eye surveillance give the best long-term results.

\section*{Prevention and Self-Care}
JIA cannot be prevented, but its effects can be minimised through early recognition and consistent, well-organised care.

\begin{itemize}
    \item See a doctor promptly for joint swelling, a persistent limp, or unexplained stiffness in a child
    \item Keep up with physiotherapy and daily exercises at home
    \item Encourage normal activity and play, and rest joints during flares before gently resuming movement
    \item Ensure the child attends every eye-screening appointment
    \item Support healthy eating, adequate sleep, and a normal school and social life
    \item Give medicines exactly as prescribed and report side effects early
    \item Seek support from the treatment team, patient organisations, and family counselling
\end{itemize}

\section*{Sources and Further Reading}
The following organisations provide reliable, evidence-based information on juvenile arthritis for parents, patients, and healthcare students.

\begin{itemize}
    \item NHS. \textit{Juvenile idiopathic arthritis (JIA).} https://www.nhs.uk/conditions/juvenile-idiopathic-arthritis/
    \item Arthritis Foundation. \textit{Juvenile arthritis information and resources.} https://www.arthritis.org/
    \item Versus Arthritis. \textit{Juvenile idiopathic arthritis for young people and families.} https://www.versusarthritis.org/about-arthritis/conditions/juvenile-idiopathic-arthritis/
    \item Mayo Clinic. \textit{Juvenile idiopathic arthritis.} https://www.mayoclinic.org/diseases-conditions/juvenile-idiopathic-arthritis/
    \item MedlinePlus. \textit{Juvenile idiopathic arthritis.} https://medlineplus.gov/juvenilerheumatoidarthritis.html
    \item National Institute of Arthritis and Musculoskeletal and Skin Diseases (NIAMS). \textit{Juvenile idiopathic arthritis.} https://www.niams.nih.gov/health-topics/juvenile-idiopathic-arthritis
\end{itemize}
